\section{System Architecture}
\label{sec:architecture}

This section presents the architecture of the proposed SysSecOps system for a hybrid
cloud environment. The design places a quantitative, multi-scanner \emph{security gate}
between the synthesis and the deployment of Infrastructure-as-Code (IaC), so that no
change is applied to either the public-cloud or the on-premises side until it has been
statically analysed, costed, checked against live compliance state, and assigned a
single risk score. The same engine governs both sides of the hybrid environment, which
gives the process a uniform and reproducible security behaviour.

\subsection{Design goals}
\label{subsec:arch-goals}

The architecture is driven by the research objectives stated in Chapter~\ref{Chapter1}
and refined into the following concrete design goals:

\begin{itemize}
    \item \textbf{Gate before deploy.} Security evaluation must occur on the
    \emph{synthesized} artifact (the CloudFormation template or the resolved Ansible
    play), not merely on the higher-level source, because the synthesized artifact is
    the last representation before any resource is provisioned.
    \item \textbf{One engine for the whole hybrid.} The AWS (public) side and the
    on-premises (private) side must be scored with identical severity weights,
    deduplication rules, and decision thresholds, so that decisions are comparable
    across the trust boundary.
    \item \textbf{Graceful degradation.} A missing scanner binary or absent cloud
    credentials must never crash the pipeline; the affected component degrades to a
    well-defined ``skipped'' status while the gate still produces a decision.
    \item \textbf{Auditability.} Every decision (\emph{pass}, \emph{review},
    \emph{reject}) is persisted as an immutable record together with the acting
    identity, enabling post-hoc review.
    \item \textbf{Feedback-driven remediation.} When code is generated, gate findings
    are injected back into the generation prompt so that unsafe code is automatically
    regenerated rather than merely blocked.
    \item \textbf{Observability.} Every deployed target is continuously monitored and
    every gate decision is published as telemetry for dashboards and automated
    remediation.
\end{itemize}

\subsection{Three-zone reference architecture}
\label{subsec:arch-zones}

The system is organised into three cooperating zones, summarised in
Table~\ref{tab:zones}. Zone~1 produces IaC; Zone~2 evaluates and gates it; Zone~3
provisions the hybrid infrastructure and continuously observes it. The risk-scoring
engine sits at the boundary of Zone~1 and Zone~2: it consumes scanner output and emits
the decision that governs whether code advances to deployment or returns to Zone~1 for
regeneration.

\begin{table}[ht]
\centering
\caption{The three zones of the proposed architecture.}
\label{tab:zones}
\begin{tabular}{|C{1.6cm}|L{3.3cm}|L{7.0cm}|}
\hline
\textbf{Zone} & \textbf{Name} & \textbf{Responsibility} \\
\hline
Zone~1 & AI + IaC local development & Generate IaC from a natural-language prompt;
local validation (\texttt{cdk synth}, \texttt{cdk diff}). \\
\hline
Zone~2 & IaC security gate & Multi-scanner analysis, cross-source deduplication, risk
scoring, and the deploy decision. \\
\hline
Zone~3 & Hybrid infrastructure and operations monitoring & Provision AWS and on-premises
resources; publish telemetry; run the drift/remediation operational loop. \\
\hline
\end{tabular}
\end{table}

Figure~\ref{fig:arch-overview} shows the three zones and the principal flow of control
between them. The solid arrows denote the forward path of a change from prompt to
deployed and monitored infrastructure; the dashed arrow denotes the remediation feedback
that returns rejected generated code to Zone~1.

\begin{figure}[ht]
\centering
\begin{tikzpicture}[node distance=6mm]
% Zone 1
\node[proc] (ui) {UI / CLI};
\node[proc, below=14mm of ui] (gen) {AI generation\\(Bedrock / OpenRouter)};
\node[io, below=14mm of gen] (app) {IaC\\(CDK / Ansible)};
\begin{scope}[on background layer]
  \node[zone, fit=(ui)(gen)(app), label=above:{\footnotesize\textbf{Generation}}] (z1) {};
\end{scope}
% Zone 2
\node[proc, right=36mm of ui] (synth) {Synthesize\\(\texttt{cdk synth})};
\node[proc, below=12mm of synth] (gate) {Multi-scanner\\gate};
\node[dec, below=12mm of gate] (dec) {score?};
\begin{scope}[on background layer]
  \node[zone, fit=(synth)(gate)(dec), label=above:{\footnotesize\textbf{Security gate}}] (z2) {};
\end{scope}
% Zone 3
\node[proc, right=24mm of synth] (deploy) {Deploy\\(AWS + on-prem)};
\node[proc, below=12mm of deploy] (mon) {Monitor\\(Phase~4 loop)};
\node[store, below=14mm of mon] (audit) {Audit\\\& metrics};
\begin{scope}[on background layer]
  \node[zone, fit=(deploy)(mon)(audit), label=above:{\footnotesize\textbf{Hybrid infra \& ops}}] (z3) {};
\end{scope}
% Edges
\draw[flow] (ui) -- (gen);
\draw[flow] (gen) -- (app);
\draw[flow] (app.east) -| (3,0) |- (synth.west) node[pos=0.7,above, font=\scriptsize]{templates};
\draw[flow] (synth) -- (gate);
\draw[flow] (gate) -- (dec);
\draw[flow] (dec.east) -| (8.6,0) |- (deploy.west) node[pos=0.75,left,font=\scriptsize,rotate=90]{pass / approved};
\draw[flow] (deploy) -- (mon);
\draw[flow] (mon) -- (audit);
\draw[flowback] (dec.west) -| (3.5,-3) |- (gen.east)
      node[pos=0.25,below,rotate=90,font=\scriptsize]{reject: inject findings};
\end{tikzpicture}
\caption{High-level three-zone architecture and end-to-end flow. Solid arrows are the
forward path; the dashed arrow is the regeneration feedback loop.}
\label{fig:arch-overview}
\end{figure}

The three zones are intentionally loosely coupled. Zone~1 knows nothing about the
internal scoring rules of Zone~2 --- it only consumes the machine-readable findings that
Zone~2 emits. Zone~3 knows nothing about how a decision was reached --- it only consumes
the decision event. This decoupling is what allows the generation backend, the set of
scanners, and the monitoring stack to evolve independently, and it is the reason the same
gate can serve both the public and the private side without modification. The remainder
of this subsection describes each zone, and the tools it uses, in turn.

\subsubsection{Zone~1 --- AI-assisted generation and local validation}
\label{subsubsec:zone1}

Zone~1 is where Infrastructure-as-Code originates. It exposes two entry surfaces --- a
command-line interface and a Streamlit operator console --- and turns a natural-language
prompt into synthesizable IaC. Code generation is delegated to a large-language-model
backend behind a provider abstraction, so that the concrete model is a configuration
choice rather than a code dependency; the reference implementation supports Amazon
Bedrock and OpenRouter. The generated artifact is an AWS Cloud Development Kit (CDK)
Python application for the public side and/or a set of Ansible playbooks for the private
side.

Before any artifact leaves Zone~1, it is validated locally so that only well-formed input
reaches the gate. For the public side, \texttt{cdk synth} synthesizes the CloudFormation
template and \texttt{cdk diff} summarises the change against the currently deployed stack;
for the private side, \texttt{ansible-playbook -{}-syntax-check} (and an optional
\texttt{-{}-check} dry run) validates the plays. Zone~1 also closes the
\emph{generate~$\rightarrow$~synth~$\rightarrow$~gate~$\rightarrow$~regenerate} feedback
loop: when the gate rejects generated code, the findings are serialised back into the
next prompt so that the model remediates its own output within a bounded attempt budget.
Table~\ref{tab:zone1-tools} lists the tools used in this zone.

\begin{table}[ht]
\centering
\caption{Tools used in Zone~1 (generation and local validation).}
\label{tab:zone1-tools}
\begin{tabular}{|L{4.6cm}|L{9.6cm}|}
\hline
\textbf{Tool} & \textbf{Role in Zone~1} \\
\hline
Amazon Bedrock / OpenRouter & Interchangeable LLM backends that generate CDK or Ansible
code from a prompt. \\
\hline
AWS CDK (Python) & Expresses public-cloud infrastructure as code and synthesizes it to
CloudFormation. \\
\hline
\texttt{cdk synth} / \texttt{cdk diff} & Produces the CloudFormation template and the
change summary used for local validation. \\
\hline
Ansible & Expresses on-premises configuration as playbooks. \\
\hline
\texttt{ansible-playbook -{}-syntax-check} & Validates playbook syntax (and, optionally,
a dry-run \texttt{-{}-check}) before gating. \\
\hline
\end{tabular}
\end{table}

\subsubsection{Zone~2 --- the multi-scanner security gate}
\label{subsubsec:zone2}

Zone~2 is the core contribution of the architecture: a quantitative security gate that
consumes the synthesized artifact and emits a single decision. It runs a set of
independent, best-of-breed scanners, normalises their heterogeneous outputs into a common
finding schema, deduplicates findings that several tools report for the same
resource/template/category, aggregates the surviving findings into a numeric risk score,
and maps that score to \emph{pass}, \emph{review}, or \emph{reject}. Each scanner is
wrapped in an adapter that never raises: a missing binary or an absent credential
degrades to a well-defined \emph{skipped} status, and the gate still produces a decision
(the graceful-degradation contract).

The tools are deliberately complementary and cover four dimensions of risk ---
\emph{misconfiguration}, \emph{cost}, \emph{live compliance}, and \emph{code-level
insecurity}. Static policy-as-code analysis (Checkov) and CloudFormation validation
(\texttt{cfn-lint}) inspect the public-side template; \texttt{ansible-lint} inspects the
private-side plays; a built-in regular-expression scanner detects hard-coded secrets;
Infracost supplies the estimated monthly cost that drives the cost dimension; AWS Config
supplies the count of live non-compliant rules; and a machine-learning code-risk model
(trained on Bandit and Semgrep features, see Section~\ref{sec:risk-scoring}) contributes a
probability of insecurity. Table~\ref{tab:zone2-tools} lists these tools and the signal
each contributes.

\begin{table}[ht]
\centering
\caption{Tools used in Zone~2 (the security gate) and the signal each contributes.}
\label{tab:zone2-tools}
\begin{tabular}{|L{3.7cm}|L{7.0cm}|L{3.1cm}|}
\hline
\textbf{Tool} & \textbf{Signal} & \textbf{Dimension} \\
\hline
AWS CDK (\texttt{cdk synth}) & Emits the CloudFormation template that is the analysed
artifact. & pre-validation \\
\hline
\texttt{cfn-lint} & CloudFormation validity and best-practice findings. & misconfig. \\
\hline
Checkov & Policy-as-code misconfiguration findings on the templates. & misconfig. \\
\hline
\texttt{ansible-lint} & Best-practice and security findings on the on-premises plays. &
misconfig. \\
\hline
Regex secret scanner (built-in) & Hard-coded credentials and secrets. & misconfig. \\
\hline
Infracost & Estimated monthly cost / cost delta of the change. & cost \\
\hline
AWS Config & Count of non-compliant live-compliance rules. & compliance \\
\hline
ML code-risk model \newline (Bandit + Semgrep $\rightarrow$ logistic regression) &
Probability that the generated code is insecure, converted to points. & code risk \\
\hline
\end{tabular}
\end{table}

\subsubsection{Zone~3 --- hybrid infrastructure and operations}
\label{subsubsec:zone3}

Zone~3 provisions the hybrid infrastructure that survived the gate and continuously
observes it. The public side is provisioned on Amazon Web Services (AWS) by deploying the
CDK-synthesized CloudFormation; the private side consists of on-premises Linux nodes
configured by Ansible. After a successful deployment, a three-layer monitoring mesh is
attached to the public resources and the on-premises nodes are registered so that both
sides are watched by a single operational loop.

On the public side, AWS Systems Manager (SSM) provides Run Command, State Manager,
Patch Manager, Inventory, and Parameter Store; Amazon EventBridge carries compliance and
stack-rollback events; Amazon CloudWatch stores the gate metrics, alarms, dashboards, and
Application Insights views; AWS CloudTrail records the API audit trail; an AWS Lambda
function implements the event-driven operational loop (drift reaction and pipeline
re-trigger); and Amazon SNS delivers operator notifications. On the private side, the SSM
Agent together with an SSM \emph{hybrid activation} registers each on-premises node as a
managed instance (\texttt{mi-*}), so the same Run Command, drift-detection, patch, and
compliance machinery used for EC2 also covers the on-premises fleet without a separate
agent stack. Table~\ref{tab:zone3-tools} summarises the tooling.

\begin{table}[ht]
\centering
\caption{Tools used in Zone~3 (hybrid infrastructure and operations).}
\label{tab:zone3-tools}
\begin{tabular}{|C{1.9cm}|L{4.0cm}|L{7.6cm}|}
\hline
\textbf{Side} & \textbf{Tool / service} & \textbf{Role} \\
\hline
Public & AWS CDK / CloudFormation & Provisions and updates the public-cloud resources. \\
\hline
Public & EC2, S3, RDS, VPC & The workload substrate that the CDK app declares. \\
\hline
Public & AWS Systems Manager & Run Command, State Manager, Patch, Inventory, and
Parameter Store (the gate-decision store). \\
\hline
Public & Amazon EventBridge & Event bus for compliance-change and stack-rollback events. \\
\hline
Public & Amazon CloudWatch & Gate metrics, alarms, dashboards, and Application Insights. \\
\hline
Public & AWS CloudTrail & Multi-region API audit trail feeding CloudWatch Logs. \\
\hline
Public & AWS Lambda & The ops-loop function: drift reaction and pipeline re-trigger. \\
\hline
Public & Amazon SNS & Operator notifications for alarms and rollbacks. \\
\hline
Private & Ansible & Configuration management of the on-premises Linux nodes. \\
\hline
Private & SSM Agent + hybrid activation & Registers each node as an \texttt{mi-*} managed
instance so it joins the same operational loop. \\
\hline
Shared & Secure connectivity link & Joins the two sides through one of several options
(see Section~\ref{subsubsec:arch-connectivity}). \\
\hline
\end{tabular}
\end{table}

\subsection{End-to-end data flow}
\label{subsec:arch-flow}

A single request flows through the system as an evolving chain of artifacts:

\begin{enumerate}
    \item A natural-language prompt (or an existing, ``bring-your-own'' project) enters
    the system through the command-line interface or the operator console.
    \item Zone~1 produces IaC: a synthesizable AWS Cloud Development Kit (CDK) Python
    application and/or a set of Ansible playbooks.
    \item The IaC is synthesized: \texttt{cdk synth} emits CloudFormation templates,
    while \texttt{ansible-playbook -{}-syntax-check} validates the playbooks.
    \item Zone~2 evaluates the synthesized artifact with several independent scanners,
    normalises their outputs into a common finding schema, deduplicates across sources,
    and computes a risk score.
    \item The decision function maps the score to \emph{pass}, \emph{review}, or
    \emph{reject}. The full gate report is persisted, and an approval or rejection
    record is written.
    \item On \emph{pass} (or an approved \emph{review}), the change is deployed; on
    \emph{reject}, generated code is optionally regenerated with the findings injected
    back into the prompt.
    \item Zone~3 records the outcome, publishes metrics and events, and attaches
    post-deployment security monitoring.
\end{enumerate}

Figure~\ref{fig:arch-dataflow} traces the same request as a chain of artifacts, making
explicit how each stage transforms its input into a more concrete representation. The
important property is that the security decision is taken on the \emph{synthesized
template}, not on the prompt or the high-level source: this is the last artifact before
provisioning and therefore the earliest point at which the exact resources, their
properties, and their relationships are fully known.

\begin{figure}[ht]
\centering
\begin{tikzpicture}[node distance=7mm and 9mm]
\node[io] (prompt) {Prompt /\\project};
\node[proc, right=of prompt] (code) {Generated\\code};
\node[proc, right=of code] (tmpl) {Synthesized\\template};
\node[proc, right=of tmpl] (find) {Normalized\\findings};
\node[dec, below=10mm of find] (score) {score \&\\decision};
\node[store, left=of score] (report) {Gate\\report};
\node[proc, left=of report] (deployN) {Deploy};
\node[store, left=of deployN] (tele) {Telemetry};
\draw[flow] (prompt) -- (code);
\draw[flow] (code) -- (tmpl);
\draw[flow] (tmpl) -- (find);
\draw[flow] (find) -- (score);
\draw[flow] (score) -- (report);
\draw[flow] (report) -- (deployN) node[pos=0.5,above,font=\scriptsize]{pass};
\draw[flow] (deployN) -- (tele);
\draw[flowback] (score.east) -| ++(6mm,4) -| (code.north)
      node[pos=0.25,below,font=\scriptsize]{reject};
\end{tikzpicture}
\caption{Artifact lifecycle of a single request. The decision is computed on the
synthesized template; a rejection feeds findings back into generation.}
\label{fig:arch-dataflow}
\end{figure}

\subsection{Hybrid deployment topology}
\label{subsec:arch-topology}

The public side is provisioned on Amazon Web Services (AWS) through CDK-synthesized
CloudFormation. The private side consists of on-premises Linux nodes managed by Ansible.
The two sides are joined by an encrypted link, and the architecture deliberately treats
\emph{how} that link is established as a pluggable concern:
Section~\ref{subsubsec:arch-connectivity} surveys the available connectivity options. The
on-premises nodes are registered as AWS Systems Manager (SSM) managed instances so that
their compliance state is visible to the same operational loop that watches the cloud
side. No public inbound administrative access is required: the control plane traverses the
encrypted link, which keeps the hybrid trust boundary small and observable.
Figure~\ref{fig:arch-topology} depicts the topology with the connectivity method shown as
an abstract secure channel.

\begin{figure}[ht]
\centering
\begin{tikzpicture}[node distance=8mm]
% AWS side
\node[proc] (vpc) {VPC / EC2 / S3 / RDS};
\node[proc, below=of vpc] (ssm) {SSM + EventBridge\\+ CloudWatch};
\begin{scope}[on background layer]
  \node[zone, fit=(vpc)(ssm), label=above:{\footnotesize\textbf{AWS (public)}}] (aws) {};
\end{scope}
% On-prem side
\node[proc, right=42mm of vpc] (node) {On-prem Linux\\node (Ansible)};
\node[proc, below=of node] (agent) {SSM agent +\\VPN client};
\begin{scope}[on background layer]
  \node[zone, fill=orange!5, fit=(node)(agent), label=above:{\footnotesize\textbf{On-premises (private)}}] (onprem) {};
\end{scope}
% Tunnel
\draw[{Stealth[length=2.4mm]}-{Stealth[length=2.4mm]}, very thick, dashed]
     (aws.east) -- (onprem.west)
     node[midway, above, font=\scriptsize, align=center]{secure\\channel};
\draw[flow] (agent.south) |- ($(ssm.south)+(0,-6mm)$) -| (ssm.south)
     node[pos=0.25,below,font=\scriptsize]{compliance events};
\end{tikzpicture}
\caption{Hybrid deployment topology: AWS resources and on-premises nodes joined by a
secure channel (any of the options in Table~\ref{tab:connectivity}), with on-premises
nodes registered as SSM managed instances.}
\label{fig:arch-topology}
\end{figure}

\subsubsection{Connectivity options between the public and private sides}
\label{subsubsec:arch-connectivity}

Joining the public and private sides is a design choice with several viable options rather
than a single fixed mechanism. The architecture treats connectivity as a pluggable
concern: the gate, the SSM registration, and the operational loop are indifferent to
\emph{how} the link is established, as long as the on-premises nodes can reach the regional
AWS endpoints and the two sides can exchange traffic securely. Table~\ref{tab:connectivity}
compares the practical options, ordered from the lightest-weight software overlay to a
dedicated physical link, so that a deployment can select the method that matches its
latency, bandwidth, cost, and compliance requirements.

\begin{table}[ht]
\centering
\caption{Options for connecting the public (AWS) and private (on-premises) sides.}
\label{tab:connectivity}
\begin{tabular}{|L{3.5cm}|L{5.6cm}|L{4.9cm}|}
\hline
\textbf{Method} & \textbf{Mechanism} & \textbf{Trade-offs} \\
\hline
Mesh overlay VPN \newline (Tailscale~/~WireGuard) &
Userspace WireGuard tunnels with automatic NAT traversal and identity-based access
control lists; each node runs a lightweight agent. & Fast to set up, no public inbound
ports, per-node identity; depends on a software agent and a coordination service. \\
\hline
AWS Site-to-Site VPN \newline (IPsec) & Managed IPsec tunnels between an on-premises
customer gateway device and an AWS virtual private gateway or transit gateway. &
Standards-based and fully encrypted over the public Internet; needs a compatible gateway
device and is bandwidth-limited per tunnel. \\
\hline
AWS Direct Connect & A dedicated private physical circuit from the on-premises data
centre into AWS, bypassing the public Internet. & Consistent low latency and high, stable
bandwidth; higher cost and provisioning lead time, and usually paired with a VPN for
encryption. \\
\hline
AWS Transit Gateway~/ \newline VPN CloudHub & A managed routing hub that interconnects
many VPCs and multiple on-premises sites in a hub-and-spoke topology. & Centralises and
scales routing across sites; adds a managed routing tier and its associated cost. \\
\hline
SSM hybrid activations \newline over public endpoints & No network tunnel: the SSM Agent
reaches the regional Systems Manager endpoints directly over TLS~443. & Simplest option
for management-plane reach only; provides no private data-plane connectivity between
workloads. \\
\hline
\end{tabular}
\end{table}

These options are not mutually exclusive: a lightweight mesh overlay can carry workload
and control traffic while an SSM hybrid activation provides the operational loop with
management-plane reach, and the heavier options (IPsec Site-to-Site VPN, Direct Connect,
or a Transit Gateway hub) can be substituted where standards-based encryption, dedicated
bandwidth, or multi-site routing is required. Crucially, the choice of connectivity is
orthogonal to the rest of the pipeline: none of these methods changes the gate, the risk
scoring, or the deployment-governance behaviour, which is what allows the architecture to
remain portable across different hybrid networking substrates.

\subsection{Component map and interfaces}
\label{subsec:arch-components}

Table~\ref{tab:modules} lists the top-level software modules and their responsibilities.
The modules communicate through stable contracts rather than shared internal state, which
keeps the generation, evaluation, orchestration, and monitoring concerns decoupled.

\begin{table}[ht]
\centering
\caption{Top-level modules and their responsibilities.}
\label{tab:modules}
\begin{tabular}{|L{5.5cm}|L{9.0cm}|}
\hline
\textbf{Module} & \textbf{Responsibility} \\
\hline
\texttt{generation/} & Zone~1 generation: provider-abstracted code generation
(Bedrock~/~OpenRouter) and the CDK regeneration loop. \\
\hline
\texttt{security\_gate/} & Zone~2 gate: template and playbook analysis, scanner adapters,
deduplication, risk scoring, and report assembly. \\
\hline
\texttt{execution/} & Subprocess execution backbone with a stable result contract. \\
\hline
\texttt{pipeline/} & Orchestration and governance: synth--gate--diff--deploy, the approval
audit trail, credentials, and the Phase~4 emitters. \\
\hline
\texttt{monitoring/} & Phase~4 observability: CloudTrail, CloudWatch metrics and
dashboards, the operational-loop Lambda, and on-premises SSM registration. \\
\hline
\texttt{scripts/} & Command-line entry points for the CDK-only and the hybrid flows. \\
\hline
\texttt{ui/} & Streamlit operator console and the multi-project registry. \\
\hline
\texttt{risk\_scoring/} & Training and evaluation code for the machine-learning
code-risk model consumed by the gate. \\
\hline
\end{tabular}
\end{table}

\subsection{Key data contracts}
\label{subsec:arch-contracts}

Three contracts hold the system together and are deliberately kept stable so that
independent consumers (the CLI, the console, the dashboards) do not break:

\begin{itemize}
    \item \textbf{Gate report.} A JSON document containing the run identifier,
    timestamp, numeric \texttt{score}, \texttt{decision}, human-readable
    \texttt{message}, a per-component \texttt{components} breakdown, the deduplicated
    \texttt{findings}, and a \texttt{scanner\_status} map that records whether each
    scanner ran, was skipped, or was unavailable.
    \item \textbf{Execution result.} Every external command returns exactly
    \texttt{\{return\_code, output\}}, with standard output and standard error merged to
    preserve event ordering for the console and the logs.
    \item \textbf{Gate-decision event.} A compact event (target name, score, decision,
    timestamp) is persisted to SSM Parameter Store and published to Amazon EventBridge,
    keyed by the target name (a stack name for CDK, or \texttt{ansible-<node>} for the
    on-premises side) so that the two sides never collide.
\end{itemize}

\subsection{Security and trust model}
\label{subsec:arch-trust}

Credentials are never passed as command-line arguments; they are resolved through the
standard provider chains and injected into subprocesses through the environment. AWS
access follows the default credential chain with a single-sign-on fallback, and
generation-provider keys are read from the environment. Deployed workloads follow the
principle of least privilege, and each privileged action is recorded in an append-only
audit trail. The uniform graceful-degradation contract guarantees that the absence of a
tool or a credential weakens visibility but never silently disables the gate.
